% ==========================================
% SECCIÓN 3: ARCHIVOS
% ==========================================
% TIP: Esta sección describe cómo se manejan los archivos en C, incluyendo la lectura y escritura de archivos, los modos de apertura, y ejemplos prácticos de uso.
%
% TIP: Las rutas de imágenes son relativas a main.tex.
%
% ==========================================
\section{Archivos en C}
\label{sec:archivos}

% ------------------------------------------
% Responsable: Cristian Delgado Lucero Isaac
% ------------------------------------------

% TODO (Cristian): Describir cómo se manejan los archivos en C, incluyendo la lectura y escritura de archivos, los modos de apertura, y ejemplos prácticos de uso.
% TODO (Cristian): Incluir ejemplos de código para cada tipo de archivo.

La evolución de la informática hacia arquitecturas de procesamiento paralelo ha exigido una redefinición constante de las abstracciones que permiten a los programas interactuar con el almacenamiento persistente. En el contexto de la programación de sistemas en lenguaje C y el desarrollo de aplicaciones paralelas, el concepto de "archivo" trasciende la mera secuencia de bytes en un disco para convertirse en un recurso compartido cuya gestión requiere una sincronización meticulosa, una comprensión de la herencia de recursos entre procesos y una optimización profunda de los mecanismos de transferencia de datos.

\subsection{Definiciones Fundamentales}

Para comprender la manipulación de datos en cómputo paralelo, es necesario definir las entidades básicas con las que interactúa el sistema operativo y el programador. Estas entidades constituyen los bloques fundamentales sobre los cuales se construyen los mecanismos de almacenamiento, comunicación y sincronización en entornos de ejecución concurrente. A continuación, se presentan dos conceptos esenciales: el archivo como unidad de almacenamiento persistente y el archivo fuente como unidad de entrada al proceso de traducción.

\subsubsection{El Concepto de Archivo}

Desde una perspectiva de abstracción de sistemas, un archivo se define como una colección de datos con nombre dentro de un sistema de archivos, estructurada esencialmente como una secuencia de bytes. Esta definición implica dos componentes críticos:

\begin{table}[H]
    \centering
    \begin{tabular}{|p{3cm}|p{10cm}|}
    \hline
    \textbf{Componente} & \textbf{Descripción} \\
    \hline
    Carga útil (Payload) & La secuencia de bytes que representa la información real (texto, imágenes, datos binarios). \\
    \hline
    Metadatos & Información sobre el archivo, como su nombre, tamaño, propietario, permisos de acceso y marcas de tiempo (creación, modificación), los cuales se almacenan en estructuras del sistema llamadas inodos. \\
    \hline
    \end{tabular}
    \caption{Componentes fundamentales de un archivo}
    \label{tab:componentes_archivo}
\end{table}

Bajo la filosofía de diseño de UNIX y POSIX, se adopta la premisa de que todo es un archivo. 
Esto significa que el sistema operativo permite a los procesos interactuar con una amplia variedad de recursos 
(discos, terminales, tuberías, sockets) utilizando la misma interfaz de llamadas al sistema 
(\texttt{open}, \texttt{read}, \texttt{write}, \texttt{close}), lo que facilita la interoperabilidad en sistemas complejos.

\subsubsection{El Archivo Fuente de C (\texttt{.c})}

Un archivo con extensión \texttt{.c} es técnicamente un archivo fuente o archivo de preprocesamiento. Sus características principales son:

\begin{table}[H]
    \centering
    \begin{tabular}{|p{3cm}|p{10cm}|}
    \hline
    \textbf{Aspecto} & \textbf{Descripción} \\
    \hline
    Contenido & Alberga el código fuente del programa, incluyendo directivas para el preprocesador (como \texttt{\#include} o \texttt{\#define}), declaraciones de tipos y definiciones de funciones. \\
    \hline
    Rol en la compilación & Es la unidad de entrada para el sistema de traducción. Cuando el preprocesador actúa sobre un archivo .c (expandiendo macros e incluyendo archivos de cabecera .h), el resultado se denomina unidad de traducción (translation unit). \\
    \hline
    Transformación & Esta unidad de traducción es lo que el compilador recibe finalmente para generar un archivo de código objeto, que posteriormente será enlazado para crear el programa ejecutable. \\
    \hline
    \end{tabular}
    \caption{Rol del archivo fuente en el proceso de compilación}
    \label{tab:archivo_fuente_compilacion}
\end{table}

\subsubsection{Marco Normativo y Evolución del Lenguaje C en la Gestión de Archivos}

La base de cualquier investigación técnica sobre programación de sistemas debe cimentarse en los estándares oficiales que rigen el comportamiento de las herramientas de desarrollo. El lenguaje C ha experimentado una transformación significativa con la publicación del estándar ISO/IEC 9899:2024, conocido como C23. Este estándar sucede al C17 y consolida prácticas que anteriormente eran extensiones de compiladores o parte de normativas regionales, elevándolas al estatus de norma internacional.

\subsection{La Abstracción del Flujo de Datos y el Sistema de Archivos}

En la arquitectura de sistemas operativos modernos, el archivo es la abstracción universal. Bajo la premisa de \enquote{todo es un archivo}, el sistema operativo permite interactuar de manera uniforme con discos duros, terminales, tuberías, sockets de red e incluso memoria física

\subsection{El Concepto de Stream en C}

El lenguaje C abstrae la complejidad de los dispositivos físicos mediante el concepto de \enquote{stream} (flujo). Un flujo es una secuencia de bytes que puede estar asociada a un archivo en disco, un teclado o una pantalla. Esta abstracción permite que las funciones de la biblioteca estándar operen de forma independiente del hardware subyacente.

Los flujos se clasifican en dos categorías principales, cuyas características se presentan en la Tabla~\ref{tab:clasificacion_flujos}.

\begin{table}[H]
    \centering
    \begin{tabular}{|p{3cm}|p{10cm}|}
    \hline
    \textbf{Tipo de Flujo} & \textbf{Descripción} \\
    \hline
    Flujos de Texto & Secuencias de caracteres organizadas en líneas. Pueden ocurrir transformaciones en los caracteres de fin de línea según el entorno (por ejemplo, de \texttt{\textbackslash n} a \texttt{\textbackslash r\textbackslash n} en Windows). \\
    \hline
    Flujos Binarios & Secuencias de bytes que no sufren transformación alguna, garantizando que los datos leídos coincidan exactamente con los almacenados. \\
    \hline
    \end{tabular}
    \caption{Clasificación de los flujos en C}
    \label{tab:clasificacion_flujos}
\end{table}

%\subsection{La Estructura FILE y la Gestión de Descriptores}

%Cada flujo está representado por un puntero a un objeto de tipo \texttt{FILE}. Aunque el programador trata este objeto como un tipo opaco, internamente el objeto \texttt{FILE} contiene información crítica gestionada por la biblioteca estándar de C (\texttt{glibc} en sistemas Linux): un descriptor de archivo (un entero que identifica el recurso ante el kernel), un búfer de memoria para optimizar las transferencias, indicadores de posición de lectura y escritura, así como banderas de estado para errores y detección de fin de archivo (EOF).

\begin{table}[H]
    \centering
    \begin{tabular}{|p{3cm}|p{7cm}|p{5cm}|}
    \hline
    \textbf{Característica} & \textbf{Flujo de Alto Nivel (FILE *)} & \textbf{Descriptor de Bajo Nivel (int)} \\
    \hline
    Origen & Estándar C (ISO C) & Estándar POSIX / Sistema Operativo \\
    \hline
    Buffering & Implementado en espacio de usuario (C Library) & Gestionado por el Kernel (Page Cache) \\
    \hline
    Portabilidad & Muy alta (funciona en cualquier compilador C) & Alta entre sistemas tipo UNIX \\
    \hline
    Funciones & \texttt{fopen}, \texttt{fprintf}, \texttt{fread}, \texttt{fclose} & \texttt{open}, \texttt{write}, \texttt{read}, \texttt{close} \\
    \hline
    Uso Ideal & E/S de texto, datos estructurados, facilidad & E/S paralela, sockets, control granular \\
    \hline
    \end{tabular}
    \caption{Comparación entre E/S de alto nivel (FILE *) y E/S de bajo nivel (descriptores)}
    \label{tab:comparacion_streams_descriptores}
\end{table}

\subsection{Operaciones de Acceso y Manipulación en el Estándar C}

El ciclo de vida de un archivo en un programa C comienza con el establecimiento de una conexión entre el flujo y el archivo físico, un proceso denominado \enquote{apertura}.

\subsubsection{Modos de Apertura y Semántica de Acceso}

La función \texttt{fopen} es la interfaz principal para este propósito. Los modos de apertura definen no solo la operación permitida (lectura, escritura o ambas), sino también el comportamiento del puntero de posición y el manejo del contenido preexistente.

\begin{table}[ht]
    \centering
    \begin{tabular}{|p{2cm}|p{3.5cm}|p{3cm}|p{5cm}|}
    \hline
    \textbf{Modo} & \textbf{Descripción} & \textbf{Posición Inicial} & \textbf{Efecto en Archivo Existente} \\
    \hline
    \texttt{r} & Lectura únicamente & Inicio del archivo & Ninguno \\
    \hline
    \texttt{w} & Escritura únicamente & Inicio del archivo & Truncamiento a cero bytes \\
    \hline
    \texttt{a} & Anexo (Append) & Fin del archivo & Los datos se añaden al final \\
    \hline
    \texttt{r+} & Lectura y Escritura & Inicio del archivo & Ninguno \\
    \hline
    \texttt{w+} & Lectura y Escritura & Inicio del archivo & Truncamiento a cero bytes \\
    \hline
    \texttt{a+} & Lectura y Anexo & Inicio (lectura) / Fin (escritura) & Los datos se añaden al final \\
    \hline
    \end{tabular}
    \caption{Modos de apertura de archivos en C}
    \label{tab:modos_apertura}
\end{table}

Es crucial notar que en modos como \texttt{a+}, la posición de lectura puede moverse libremente mediante funciones de posicionamiento, pero todas las operaciones de escritura se realizarán forzosamente al final del archivo, garantizando la integridad de los datos añadidos en entornos donde múltiples procesos podrían estar operando de forma secuencial.

\subsubsection{Funciones de Entrada y Salida Directa}

Para la práctica de cómputo paralelo, donde a menudo se extraen grandes bloques de información para procesarlos, las funciones \texttt{fread()} y \texttt{fwrite()} ofrecen el mejor rendimiento al evitar la sobrecarga de la interpretación de cadenas de formato propia de \texttt{fprintf()} o \texttt{fscanf()}.

La función \texttt{fread(void *ptr, size\_t size, size\_t nmemb, FILE *stream)} lee \texttt{nmemb} elementos, cada uno de tamaño \texttt{size}, desde el flujo al área de memoria apuntada por \texttt{ptr}. Por su parte, \texttt{fwrite(const void *ptr, size\_t size, size\_t nmemb, FILE *stream)} escribe bloques binarios desde la memoria al archivo, devolviendo el número de elementos escritos exitosamente, lo que permite detectar errores de espacio en disco o interrupciones en el sistema.

\subsubsection{Control de Posicionamiento}

El acceso aleatorio (random access) es una característica esencial de los archivos en disco que permiten operaciones en cualquier posición. Las funciones \texttt{fseek()} y \texttt{ftell()} (o sus variantes de mayor precisión \texttt{fseeko()} y \texttt{ftello()} en POSIX) permiten manipular el indicador de posición del flujo.

La función \texttt{fseek()} toma un desplazamiento (offset) y un punto de referencia, los cuales pueden ser:
\begin{itemize}
    \item \texttt{SEEK\_SET}: Inicio del archivo.
    \item \texttt{SEEK\_CUR}: Posición actual del indicador.
    \item \texttt{SEEK\_END}: Fin del archivo.
\end{itemize}

En el cómputo paralelo, el posicionamiento preciso permite que diferentes procesos operen en diferentes segmentos del mismo archivo, una técnica conocida como particionamiento de E/S.

\subsection{Arquitectura del Sistema de Archivos: El Rol del Inodo y los Metadatos}

Para diseñar un sistema que extraiga información de manera eficiente, es necesario comprender la estructura física y lógica en la que reside el archivo. Los sistemas operativos modernos organizan el almacenamiento basándose en el concepto de \enquote{inodo} (Index Node).

\subsubsection{Anatomía de un Inodo}

Un inodo es una estructura de datos en el disco que describe un objeto del sistema de archivos, como un archivo o un directorio. Contiene toda la información sobre el archivo excepto su nombre y los datos reales. Los metadatos almacenados en el inodo se detallan en la Tabla~\ref{tab:metadatos_inodo}.

\begin{table}[ht]
    \centering
    \begin{tabular}{|p{3.5cm}|p{9.5cm}|}
    \hline
    \textbf{Metadato} & \textbf{Descripción} \\
    \hline
    File Mode & Tipo de archivo (regular, directorio, pipe) y permisos (rwx) \\
    \hline
    Owner Info & ID de usuario (UID) y ID de grupo (GID) \\
    \hline
    File Size & Tamaño actual del archivo en bytes \\
    \hline
    Time Stamps & Creación (ctime), modificación (mtime) y acceso (atime) \\
    \hline
    Link Count & Número de nombres de archivo que apuntan a este inodo \\
    \hline
    Data Blocks & Punteros a los bloques físicos en el disco donde residen los datos \\
    \hline
    \end{tabular}
    \caption{Metadatos almacenados en un inodo}
    \label{tab:metadatos_inodo}
\end{table}

\subsubsection{Jerarquía de Directorios y Resolución de Nombres}

Un directorio es, en esencia, un archivo especial cuyo contenido es una tabla de pares (Nombre, Número de Inodo). Cuando un programa intenta abrir \enquote{\texttt{/home/usuario/datos.txt}}, el sistema operativo realiza una resolución de ruta que sigue los siguientes pasos:
\begin{enumerate}
    \item Busca el inodo del directorio raíz (\texttt{/}).
    \item Lee el contenido de \texttt{/} para encontrar el inodo de \texttt{home}.
    \item Lee el contenido de \texttt{home} para encontrar el inodo de \texttt{usuario}.
    \item Finalmente, localiza el inodo de \texttt{datos.txt} para acceder a sus bloques de datos.
\end{enumerate}

Para optimizar este proceso, el kernel mantiene una caché de entradas de directorio (\texttt{dentry cache}) en RAM, lo que reduce drásticamente las lecturas a disco necesarias para encontrar archivos frecuentemente accedidos.

\subsection{Implementación de Archivos en Memoria: \texttt{memfd\_create()}}

En sistemas de cómputo paralelo, a menudo es ineficiente escribir resultados intermedios en un disco físico para que otro proceso los lea. Linux ofrece la llamada al sistema \texttt{memfd\_create()}, que permite crear archivos anónimos que residen completamente en la memoria RAM (\texttt{tmpfs}).

Estos \enquote{archivos de memoria} se comportan como archivos regulares (permiten \texttt{read}, \texttt{write}, \texttt{lseek}, \texttt{mmap}), pero no tienen una entrada en el sistema de archivos físico (no aparecen en \texttt{ls}). Devuelven un descriptor de archivo que puede ser compartido entre procesos mediante \texttt{fork()} o enviado a través de sockets de dominio UNIX.

Adicionalmente, estos archivos soportan un mecanismo de \enquote{sellado} (sealing) mediante \texttt{fcntl()}, lo que permite que un proceso escriba datos en el archivo y luego le ponga un sello de solo lectura antes de entregarlo a otros procesos, garantizando que nadie pueda modificar los resultados originales.